\documentclass{article}

% =================================================================
% PACKAGE SETUP
% =================================================================
\usepackage{microtype}
\usepackage{graphicx}
\usepackage{subcaption}
\usepackage{booktabs} % professional-quality tables
\usepackage{multirow}
\usepackage{amsmath}
\usepackage{amssymb}
\usepackage{mathtools}
\usepackage{amsthm}
\usepackage{bm}
\usepackage{xcolor}
\usepackage{algorithm}
\usepackage{algorithmic}
\usepackage[utf8]{inputenc}
\usepackage{enumitem} % For compact lists

\usepackage[colorlinks]{hyperref}

% Use the ICML style file (proxy for 2026)
\usepackage{icml2026} 

% =================================================================
% THEOREMS & MACROS
% =================================================================
\theoremstyle{plain}
\newtheorem{theorem}{Theorem}[section]
\newtheorem{lemma}[theorem]{Lemma}
\newtheorem{proposition}[theorem]{Proposition}
\newtheorem{corollary}[theorem]{Corollary}

\theoremstyle{definition}
\newtheorem{definition}[theorem]{Definition}
\newtheorem{assumption}[theorem]{Assumption}
\newtheorem{example}[theorem]{Example}

\theoremstyle{remark}
\newtheorem{remark}[theorem]{Remark}

% Math Macros
\newcommand{\R}{\mathbb{R}}
\newcommand{\E}{\mathbb{E}}
\newcommand{\norm}[1]{\left\| #1 \right\|}
\newcommand{\inner}[2]{\langle #1, #2 \rangle}
\newcommand{\OmegaSpace}{\Omega} % Token space
\newcommand{\Pspace}{\mathcal{P}(\OmegaSpace)} % Probability measure space
\newcommand{\Pm}[1]{\mathcal{P}_{#1}(\OmegaSpace)} % Discrete measure space
\newcommand{\muN}{\mu_N} % Full context
\newcommand{\tmu}{\tilde{\mu}} % Compressed context
\newcommand{\Wass}[1]{W_{#1}} % Wasserstein distance
\newcommand{\Attn}{\text{Attn}}
\newcommand{\LambdaOp}{\Lambda_{\theta}} % Transformer Operator
\newcommand{\kernel}{\kappa}
\newcommand{\valmap}{\phi}
\newcommand{\Compress}{\mathcal{C}_m} % Compression operator

% Short title
\icmltitlerunning{Transformers as Measure Learners}

\begin{document}

\twocolumn[
% =================================================================
% TITLE BLOCK: More Theoretical and Grand
% =================================================================
\icmltitle{Transformers as Measure Learners: \\
Wasserstein Scaling Laws for KV Cache Compression}

\icmlsetsymbol{equal}{*}

\begin{icmlauthorlist}
\icmlauthor{Author One}{equal,inst1}
\icmlauthor{Author Two}{equal,inst1}
\icmlauthor{Author Three}{inst2}
\end{icmlauthorlist}

\icmlaffiliation{inst1}{Department of Computer Science, University X}
\icmlaffiliation{inst2}{Department of Mathematics, University Y}

\icmlcorrespondingauthor{Author One}{email@domain.edu}

\icmlkeywords{In-Context Learning, Long-Context Transformers, KV Cache Compression, Measure Theory, Optimal Transport, Wasserstein Geometry}

]
\printAffiliationsAndNotice{\icmlEqualContribution} 
\vskip 0.3in

\begin{abstract}
In long-context In-Context Learning (ICL), the main bottleneck is the Key--Value (KV) cache, whose memory grows linearly with sequence length. Existing compression methods prune tokens or quantize values, but treat the context as a scored sequence and offer little guidance on what \emph{optimal} compression should be.

We instead view the context as an empirical probability measure and the KV cache as its finite-support approximation. Building on universality results for Transformers on measures, we show that ICL predictions are Lipschitz with respect to the Wasserstein-1 distance between context measures, which yields a geometric error bound for KV compression. Combining this with quantization theory gives a minimax-optimal $m^{-1/d}$ scaling of ICL error in the KV budget $m$, where $d$ is the intrinsic feature dimension.

Guided by this view, we propose \textbf{Measure-Aware Compression (MAC)}, which approximates Wasserstein minimization via weighted K-Means at inference time and Sinkhorn-regularized induced attention during training. On synthetic ICL tasks and the LUNA16 medical benchmark, MAC achieves better memory--accuracy tradeoffs than strong KV baselines such as H2O and Quest.
\end{abstract}


% =================================================================
% 1. INTRODUCTION
% =================================================================
\section{Introduction}
\label{sec:intro}

Large Transformer models have emerged as powerful \emph{in-context learners}: given a long context of input--output examples, they can synthesize a predictor for new queries without updating parameters \citep{brown2020language,garg2022what}. With recent progress in long-context architectures such as LongRoPE, Unlimiformer, and proprietary systems like Gemini 1.5 and Llama 2 \citep{ding2024longrope,bertsch2024unlimiformer,team2024gemini,touvron2023llama}, context windows have reached millions of tokens, enabling foundation models for domains like medicine to reason over entire patient trajectories, imaging studies, or clinical records \citep{moor2023foundation}.

\paragraph{The KV Memory Wall of Long-Context ICL.}
Despite architectural advances, long-context In-Context Learning (ICL) on Transformers is fundamentally constrained by the \emph{Key--Value (KV) cache}. In standard self-attention \citep{vaswani2017attention}, the KV cache grows linearly $O(N)$ with sequence length $N$, and attention computation scales as $O(N^2)$. For million-token contexts, storing KV states in FP16 can require hundreds of gigabytes of memory, exceeding typical GPU budgets and making inference prohibitively slow—especially in resource-constrained environments such as hospital edge servers or autonomous labs. Alternative sequence models, including linear/streaming attention \citep{katharopoulos2020transformers} and state-space models like Mamba \citep{gu2024mamba}, reduce asymptotic complexity, but empirical evidence suggests that standard Transformers remain highly competitive ICL systems due to their flexible associative memory \citep{bietti2024birth}. This motivates a principled approach to \emph{compressing} the KV cache while preserving ICL performance.

\paragraph{Heuristic KV Compression and Its Limits.}
A rapidly growing literature attacks the KV bottleneck with engineering techniques. Token-pruning methods such as H2O, Scissorhands, SnapKV, and PyramidKV \citep{zhang2024h2o,liu2024scissorhands,li2024snapkv,zhang2024pyramidkv,ge2024model} retain “important’’ tokens based on accumulated attention scores or layer-wise saliency and evict the rest. Retrieval-based approaches, including Quest, Infini-attention, RetrievalAttention, and Unlimiformer \citep{tang2024quest,munkhdalai2024leave,gao2024retrieval,bertsch2024unlimiformer}, maintain a larger external memory and fetch relevant blocks on demand. Quantization and distillation-based methods like KIVI, attention sinks, and KV distillation \citep{liu2024kivi,xiao2024efficient,nawrot2024attention} instead compress the representation of each KV pair to a few bits or a small distilled cache.

All of these methods, however, share two structural limitations. First, they treat the context as an \emph{ordered sequence with local scores} (per-token or per-block importance) rather than as a global geometric object. Second, they lack \emph{distribution-level guarantees}: there is no theory explaining what an \emph{optimal} compressed KV cache should look like, or how the induced compression error scales with the memory budget and the structure of the underlying data.

\paragraph{Our Perspective: Transformers as Measure Learners.}
We propose a different view. Instead of thinking of the context as a scored sequence, we interpret it as an \emph{empirical probability measure} on a feature space. A sequence of token features $Z_N = \{z_i\}_{i=1}^N$ induces an empirical measure $\mu_N = \tfrac{1}{N}\sum_i \delta_{z_i}$ on a compact domain $\Omega \subset \mathbb{R}^d$. The KV cache is then a finite-support approximation $\tilde{\mu}$ with at most $m \ll N$ atoms. Under this view, KV compression is no longer a heuristic selection problem, but a \emph{measure quantization} problem: construct a sparse measure that best approximates $\mu_N$ under a suitable metric.

This perspective is supported by recent theory. Transformers have been shown to be universal approximators of continuous functionals on spaces of probability measures \citep{furuya2025transformers}, and kernel mean embeddings provide an injective representation of distributions in reproducing kernel Hilbert spaces \citep{muandet2017kernel,sriperumbudur2010hilbert}. At the same time, optimal transport (OT) and Wasserstein geometry offer natural metrics on probability measures \citep{villani2009optimal,cuturi2013sinkhorn}, and classical quantization theory characterizes how well a measure on a $d$-dimensional manifold can be approximated by $m$ atoms, with minimax error scaling as $m^{-1/d}$ \citep{graf2000foundations}. We combine these ideas to ask:

\begin{quote}
\emph{If Transformers are measure learners, what is the optimal way to compress the context measure under a fixed KV budget, and how does the induced ICL error scale with the budget and intrinsic dimension of the data?}
\end{quote}

\paragraph{From Wasserstein Stability to Optimal KV Compression.}
Our key technical observation is that attention can be viewed as an integral operator over the context measure. Building on smoothness analyses of attention \citep{castin2024how,sander2022sinkformers} and kernel embedding theory \citep{muandet2017kernel,sriperumbudur2010hilbert}, we show that under mild regularity assumptions, Transformer ICL predictions are \emph{Lipschitz} with respect to the Wasserstein-1 distance between context measures. This yields a geometric error bound: if a compressor $\mathcal{C}_m$ produces a sparse measure $\tilde{\mu}_m$ that is close to $\mu_N$ in Wasserstein-1, then the ICL prediction error is controlled by $\mathcal{W}_1(\mu_N,\tilde{\mu}_m)$. Combining this Lipschitz bound with classical quantization results \citep{graf2000foundations,villani2009optimal} yields a minimax-optimal $m^{-1/d}$ scaling law for ICL error in the KV budget $m$, where $d$ is the intrinsic feature dimension.

\paragraph{Contributions.}
We summarize our contributions as follows:
\begin{itemize}[leftmargin=*,noitemsep,topsep=2pt]
    \item \textbf{Theory: Transformers as Wasserstein-stable measure learners.} We formalize ICL as operator learning on the space of probability measures and, leveraging universality results \citep{furuya2025transformers}, show that Transformers with kernel attention approximate continuous functionals on $(\mathcal{P}(\Omega),\mathcal{W}_1)$. We prove that attention heads are Lipschitz with respect to the Wasserstein-1 distance, and derive a network-level error bound that directly links KV compression to ICL prediction error.
    \item \textbf{Optimal scaling laws for KV compression.} By coupling our Wasserstein stability bound with quantization theory \citep{graf2000foundations}, we show that if token features lie on a $d$-dimensional manifold, there exists a KV compression strategy such that the ICL error decays as $\tilde{\mathcal{O}}(m^{-1/d})$ in the KV budget $m$, and that this rate is minimax-optimal. This reveals an intrinsic-dimension bottleneck: structured scientific data (low $d$) is highly compressible, whereas high-entropy text is fundamentally harder to summarize.
    \item \textbf{Measure-Aware Compression (MAC).} Guided by this geometry, we introduce \emph{Measure-Aware Compression} (MAC), a family of practical algorithms that explicitly respect the measure structure of the context. MAC-I performs post-hoc KV compression at inference time via weighted K-Means, approximating Wasserstein minimization, while MAC-T integrates Sinkhorn-regularized induced attention \citep{cuturi2013sinkhorn,lee2019set} into training to learn compressive representations end-to-end.
    \item \textbf{Empirical validation on synthetic and medical benchmarks.} On synthetic linear regression ICL tasks, MAC follows the predicted $m^{-1/d}$ scaling, whereas score-based pruning saturates early. On the LUNA16 pulmonary nodule detection benchmark \citep{setio2017validation}, MAC achieves strong memory--accuracy tradeoffs, outperforming state-of-the-art KV compression baselines such as H2O and Quest \citep{zhang2024h2o,tang2024quest} under the same memory budget. Qualitative analysis shows that MAC preserves rare but clinically critical “outlier’’ slices in feature space that heuristic methods tend to discard.
\end{itemize}

Overall, our results suggest that Transformers are not only sequence models, but \emph{measure learners}, and that KV compression is best understood as constructing an optimal geometric summary of a context distribution rather than simply dropping low-score tokens.

% =================================================================
% 2. RELATED WORK
% =================================================================
\section{Related Work}
\label{sec:related}

We organize related work into four areas: (i) theoretical views of In-Context Learning (ICL),  
(ii) KV cache compression and efficient long-context Transformers, (iii) optimal transport and  
kernel mean embeddings, and (iv) long-context applications in medical AI.

\paragraph{ICL Theory and Transformers as Operators.}
A growing line of work studies what and how Transformers can learn in context.  
\citet{garg2022what} analyze which simple function classes can be implemented via ICL,  
while \citet{akyurek2022learning,von2023transformers,dai2022why} show that ICL can  
realize variants of gradient descent or meta-optimization, and  
\citet{ahn2024transformers,mahankali2024one} characterize linear self-attention as an  
optimal in-context learner under suitable assumptions.  
Mechanistic studies such as \citet{olsson2022context} further reveal “induction heads’’  
that implement pattern completion in context.  
Most relevant to our perspective, \citet{furuya2025transformers} prove that Transformers  
are universal approximators for functionals on probability measures, suggesting that  
they can act directly on distributions rather than purely on ordered sequences.  
However, these works assume access to the full (uncompressed) context and do not  
address memory constraints.  
In contrast, we build on the measure-theoretic view of \citet{furuya2025transformers}  
and combine it with quantization theory to obtain Wasserstein-based stability results  
and minimax error scaling laws under KV compression.

\paragraph{KV Cache Compression and Efficient Long-Context Transformers.}
The quadratic cost of self-attention \citep{vaswani2017attention} and the linear growth  
of the KV cache motivate a large literature on efficient Transformers.  
Linear or streaming attention methods such as \citet{katharopoulos2020transformers}  
and state-space models like Mamba \citep{gu2024mamba} reduce asymptotic complexity,  
yet empirical evidence suggests that standard Transformers remain highly competitive  
for ICL and long-context reasoning \citep{bietti2024birth}.  
At the system level, architectures like LongRoPE, Unlimiformer, Gemini 1.5, and  
Llama 2 \citep{ding2024longrope,bertsch2024unlimiformer,team2024gemini,touvron2023llama}  
push context lengths to millions of tokens, exacerbating the KV memory bottleneck.

KV cache compression methods can be roughly grouped into three categories.  
(1) \emph{Token pruning}: H2O, Scissorhands, SnapKV, PyramidKV, and hybrid eviction schemes  
\citep{zhang2024h2o,liu2024scissorhands,li2024snapkv,zhang2024pyramidkv,ge2024model}  
retain “important’’ tokens according to accumulated attention scores or learned saliency,  
evicting the rest.  
(2) \emph{Retrieval-based memory}: Quest, Infini-attention, RetrievalAttention, and  
related approaches \citep{tang2024quest,munkhdalai2024leave,gao2024retrieval,bertsch2024unlimiformer}  
store a larger external memory and retrieve relevant blocks on demand, effectively  
trading KV size for retrieval latency and system complexity.  
(3) \emph{Quantization and distillation}: KIVI, attention sinks, and distilled KV caches  
\citep{liu2024kivi,xiao2024efficient,nawrot2024attention} compress each KV pair to a few  
bits or a smaller distilled representation without reducing the number of tokens.  
All of these methods operate primarily at the level of individual tokens or blocks and  
are guided by local importance scores.  They do not model the context as a probability  
measure, and they provide no distribution-level guarantees that relate compression to  
ICL prediction error.  Our work instead treats the KV cache as a finite-support  
approximation of an underlying context measure and casts KV compression as a measure  
quantization problem, which leads directly to geometric error bounds and intrinsic  
scaling laws.

\paragraph{Optimal Transport, Kernel Mean Embeddings, and Attention Geometry.}
Our analysis draws on tools from optimal transport (OT), kernel methods, and recent  
smoothness studies of attention.  The Wasserstein distance and its computational  
approximations \citep{villani2009optimal,cuturi2013sinkhorn} provide natural metrics  
on probability measures and have been widely applied in modern machine learning.  
Classical quantization theory \citep{graf2000foundations} characterizes how well a  
measure supported on a $d$-dimensional manifold can be approximated by a finite-support  
measure with $m$ atoms, showing minimax rates of order $m^{-1/d}$.  
We leverage these results to derive intrinsic-dimension–dependent scaling laws for  
KV compression.  
Kernel mean embeddings \citep{muandet2017kernel,sriperumbudur2010hilbert} map probability  
measures into reproducing kernel Hilbert spaces and yield injective representations  
under characteristic kernels; this aligns naturally with viewing attention as an  
integral operator, and connects to Performer-style kernel approximations via random  
features \citep{song2020rethinking,rahimi2007random}.  
The smoothness and Lipschitz properties of attention maps have been studied by  
\citet{castin2024how} and in doubly-stochastic attention variants such as Sinkformers  
\citep{sander2022sinkformers}, while \citet{chizat2018global} analyze the dynamics of  
over-parameterized residual networks through an OT lens.  
Our work combines these ideas to show that attention is Lipschitz with respect to  
Wasserstein-1 and that KV compression error propagates through the network in a  
controlled way, yielding explicit Wasserstein-based error bounds.

\paragraph{Long-Context Modeling in Medical AI.}
Long-context reasoning is particularly important in medicine, where patient trajectories,  
multi-slice imaging studies, and multi-modal records naturally form long sequences.  
\citet{moor2023foundation} highlight the need for generalist medical foundation models  
that can integrate such rich histories under realistic computational budgets.  
LUNA16 \citep{setio2017validation} provides a benchmark for pulmonary nodule detection  
from CT scans, where each study comprises hundreds of slices; this offers a natural  
testbed for evaluating long-context models and KV compression strategies in a  
clinically meaningful setting.  
By applying our measure-aware KV compression to LUNA16 and comparing against  
state-of-the-art baselines such as H2O and Quest \citep{zhang2024h2o,tang2024quest},  
we demonstrate that geometric, distribution-level compression can preserve rare but  
clinically critical slices that heuristic pruning or retrieval tends to discard.

\paragraph{Summary.}
In summary, prior work has developed powerful operator-theoretic views of ICL and a rich  
toolkit of engineering solutions for KV compression, but lacks a unifying measure-based  
theory that explains what \emph{optimal} KV compression should look like and how its  
error scales with the memory budget and intrinsic data dimension.  
Our contribution is to interpret Transformers as \emph{measure learners}, to formalize  
KV compression as measure quantization under Wasserstein geometry, and to instantiate  
this view in practical algorithms that achieve state-of-the-art memory--accuracy  
tradeoffs on synthetic and medical long-context benchmarks.


% =================================================================
% 3. PRELIMINARIES
% =================================================================
\section{Preliminaries}
\label{sec:prelim}

We briefly summarize the measure-theoretic and optimal transport notation used in our analysis, and formalize attention as an integral operator on probability measures.

\paragraph{Context as an Empirical Measure.}
We assume token features lie in a compact domain $\OmegaSpace \subset \R^d$.  
A length-$N$ context $Z_N = \{z_i\}_{i=1}^N$ (e.g., keys, values, or hidden states) induces an empirical probability measure
\begin{equation}
    \muN \;=\; \frac{1}{N} \sum_{i=1}^N \delta_{z_i} \;\in\; \Pspace,
\end{equation}
where $\delta_{z_i}$ is the Dirac mass at $z_i$ and $\Pspace$ denotes the set of Borel probability measures on $\OmegaSpace$.  
Under this view, the KV cache is a finite-support approximation of $\muN$ with at most $m \ll N$ atoms; KV compression corresponds to replacing $\muN$ by a sparse measure $\tmu_m$.

\paragraph{Wasserstein Distance and Quantization.}
For $p \ge 1$, the $p$-Wasserstein distance between $\mu, \nu \in \Pspace$ is defined as \citep{villani2009optimal}
\begin{equation}
    \Wass{p}(\mu, \nu)
    \;\coloneqq\;
    \Bigg(
    \inf_{\pi \in \Pi(\mu, \nu)}
    \int_{\OmegaSpace \times \OmegaSpace}
    \|x - y\|^p \, d\pi(x, y)
    \Bigg)^{\!1/p},
\end{equation}
where $\Pi(\mu,\nu)$ is the set of couplings with marginals $\mu$ and $\nu$.  
We will focus on $\Wass{1}$, which metrizes weak convergence on $\Pspace$ and admits a dual formulation that directly relates Lipschitz test functions to measure perturbations \citep{villani2009optimal}.  

Given a budget $m$, a \emph{quantization} (or compression) operator produces a discrete measure
\begin{equation}
    \Compress(\muN) \;=\; \tmu_m \;=\; \sum_{j=1}^{m} w_j \, \delta_{\tilde z_j},
    \qquad \sum_{j=1}^m w_j = 1,
\end{equation}
with at most $m$ atoms.  Classical quantization theory characterizes how small $\Wass{1}(\muN,\tmu_m)$ can be as a function of $m$ and the intrinsic dimension of the support of $\muN$ \citep{graf2000foundations}; we will leverage these results in our scaling-law analysis.

\paragraph{Kernel Mean Embeddings and Attention as Integration.}
Let $k : \OmegaSpace \times \OmegaSpace \to \R$ be a positive-definite kernel with associated reproducing kernel Hilbert space (RKHS) $\mathcal{H}_K$ \citep{muandet2017kernel}.  
The \emph{kernel mean embedding} of a measure $\mu \in \Pspace$ is
\begin{equation}
    \phi(\mu) \;\coloneqq\; \int_{\OmegaSpace} k(\cdot, z)\, d\mu(z) \;\in\; \mathcal{H}_K.
\end{equation}
For characteristic kernels, this map is injective \citep{sriperumbudur2010hilbert}, providing an implicit representation of $\mu$ as a point in $\mathcal{H}_K$.

We model a single attention head as an integral operator over the context measure.  
Given a query $q \in \R^{d_q}$, a kernel family $\kappa(q,\cdot)$, and a value map $\valmap : \OmegaSpace \to \R^{d_v}$, we define
\begin{equation}
    \label{eq:integral_attn}
    \Attn(q; \mu)
    \;\coloneqq\;
    \frac{
        \int_{\OmegaSpace} \kappa(q, z)\,\valmap(z)\, d\mu(z)
    }{
        \int_{\OmegaSpace} \kappa(q, z)\, d\mu(z)
    }.
\end{equation}
Standard softmax attention \citep{vaswani2017attention} is recovered by choosing
$\kappa(q,z) = \exp(q^\top k)$ and letting $z$ bundle keys and values $(k,v)$ with $\valmap(z)=v$.  
In the discrete case $\mu = \muN$, the integrals in \eqref{eq:integral_attn} reduce to finite sums over tokens; when $\mu$ is compressed to $\tmu_m$, the same operator is applied to a sparse measure with $m$ atoms.  

This measure-based formulation makes the dependence of attention on $\mu$ explicit, and will allow us to study the Lipschitz continuity of $\Attn(\cdot;\mu)$ with respect to $\Wass{1}$ and to propagate KV compression error through multi-layer Transformers in Section~\ref{sec:theory}.

% =================================================================
% 4. THEORETICAL FRAMEWORK
% =================================================================
\section{Theoretical Framework}
\label{sec:theory}

We now formalize Transformers as \emph{measure learners} and derive the key ingredients of our framework:  
(i) universality on spaces of probability measures, (ii) Wasserstein stability of attention, and  
(iii) minimax scaling laws for KV compression.

Throughout, $\OmegaSpace \subset \R^d$ is compact, $\Pspace$ is the set of Borel probability measures on $\OmegaSpace$,  
and $\Wass{1}$ denotes the 1-Wasserstein distance on $\Pspace$ \citep{villani2009optimal}.  
A context $Z_N = \{z_i\}_{i=1}^N$ induces an empirical measure $\muN$, and a compressed KV cache corresponds to a sparse measure $\tmu_m$ with at most $m$ atoms.

\subsection{Setup: ICL as Operator Learning on Measures}
\label{sec:theory_setup}

We model In-Context Learning as \emph{operator learning} on measures.  
Given a context measure $\mu \in \Pspace$ and a query $q \in \OmegaSpace$, the ideal in-context predictor is a functional
\begin{equation}
    F: \Pspace \times \OmegaSpace \to \mathcal{Y}, 
    \qquad (\mu,q) \mapsto F(\mu,q),
\end{equation}
where $\mathcal{Y}$ is the output space (e.g., logits, regression targets).  
Transformers with kernel attention (Eq.~\eqref{eq:integral_attn}) define a parametric family of such functionals, which we denote by $\LambdaOp_\theta(\mu,q)$.

We make a mild regularity assumption on the feature space and kernel family.

\begin{assumption}[Regularity]
\label{ass:regularity}
\begin{enumerate}[leftmargin=*,nosep]
    \item $\OmegaSpace \subset \R^d$ is compact.
    \item The kernel family $\{\kappa(q,\cdot) : q \in \OmegaSpace\}$ is characteristic and Lipschitz, and the value map $\valmap$ is bounded and Lipschitz.
    \item For all relevant $\mu \in \Pspace$ and $q \in \OmegaSpace$, the denominator in Eq.~\eqref{eq:integral_attn} is bounded below: 
    $\int \kappa(q,z)\, d\mu(z) \ge \delta > 0$.
\end{enumerate}
\end{assumption}

Assumption~\ref{ass:regularity} is standard in kernel mean embedding and OT analyses  
\citep{muandet2017kernel,sriperumbudur2010hilbert,villani2009optimal,castin2024how} and holds for common choices such as Gaussian kernels on bounded domains.

\subsection{Universality on Spaces of Measures}
\label{sec:theory_universality}

Recent work shows that Transformers are universal approximators for functionals on probability measures \citep{furuya2025transformers}.  
We adapt this result to our kernel-attention formulation.

\begin{theorem}[Universal Approximation on Measures]
\label{thm:universality}
Let $C(\Pspace \times \OmegaSpace)$ denote the space of continuous functionals on $\Pspace \times \OmegaSpace$,  
where $\Pspace$ is equipped with the weak topology metrized by $\Wass{1}$.  
Under Assumption~\ref{ass:regularity}, for any $F \in C(\Pspace \times \OmegaSpace)$ and any $\varepsilon > 0$,  
there exists a Transformer $\LambdaOp_\theta$ with finitely many kernel-attention heads and MLP layers such that
\begin{equation}
    \sup_{\mu \in \Pspace,\, q \in \OmegaSpace}
    \big\| F(\mu,q) - \LambdaOp_\theta(\mu,q) \big\|
    < \varepsilon.
\end{equation}
\end{theorem}

The proof (Appendix~\ref{app:proof_universality}) combines the universality of Transformers on measures \citep{furuya2025transformers}  
with the injectivity of kernel mean embeddings under characteristic kernels \citep{muandet2017kernel,sriperumbudur2010hilbert}.  
Intuitively, attention heads generate a rich algebra of measure-dependent features that separates points in $\Pspace \times \OmegaSpace$,  
and MLPs provide the final nonlinear mixing.

\subsection{Wasserstein Stability of Attention}
\label{sec:theory_lipschitz}

To understand KV compression, we need to quantify how sensitive attention is to perturbations of the context measure.  
We view a single attention head as the map $\mu \mapsto \Attn(q;\mu)$ defined in Eq.~\eqref{eq:integral_attn}.

\begin{lemma}[Attention is Lipschitz in $\Wass{1}$]
\label{lemma:lipschitz}
Under Assumption~\ref{ass:regularity}, for any fixed query $q \in \OmegaSpace$,  
the attention operator $\mu \mapsto \Attn(q;\mu)$ is Lipschitz with respect to $\Wass{1}$:
\begin{equation}
    \big\| \Attn(q;\mu) - \Attn(q;\nu) \big\|
    \;\le\;
    L_{\Attn} \, \Wass{1}(\mu,\nu),
    \quad \forall \mu,\nu \in \Pspace,
\end{equation}
where $L_{\Attn}$ depends on the Lipschitz constants and bounds of $\kappa$ and $\valmap$,  
the diameter of $\OmegaSpace$, and the lower bound $\delta$ on the denominator.
\end{lemma}

The proof (Appendix~\ref{app:proof_lipschitz}) uses the Kantorovich–Rubinstein duality for $\Wass{1}$ and a quotient-rule argument:  
both the numerator and denominator in Eq.~\eqref{eq:integral_attn} are Lipschitz linear functionals of $\mu$,  
and the denominator is bounded away from zero, so their ratio is Lipschitz.  
This is consistent with recent smoothness analyses of attention maps \citep{castin2024how,sander2022sinkformers}.

\subsection{Network-Level KV Compression Bound}
\label{sec:theory_compression}

KV compression replaces the full context measure $\muN$ by a sparse approximation $\tmu_m = \Compress(\muN)$ with at most $m$ atoms.  
We now show that, under mild Lipschitz assumptions on the layers, the induced ICL error is controlled by $\Wass{1}(\muN,\tmu_m)$.

Let $\LambdaOp_\theta = \Lambda^{(L)} \circ \cdots \circ \Lambda^{(1)}$ denote an $L$-layer Transformer.  
Assume each layer $\Lambda^{(l)}$ is $L_{\text{layer}}^{(l)}$-Lipschitz in its measure argument (this holds for attention + MLP + residual blocks under Assumption~\ref{ass:regularity} and standard normalization).

\begin{theorem}[KV Compression Error Bound]
\label{thm:bound}
Under Assumption~\ref{ass:regularity}, let $\muN$ be the empirical context measure and $\tmu_m = \Compress(\muN)$ any compressed measure with at most $m$ atoms. Then
\begin{equation}
    \big\|
        \LambdaOp_\theta(\muN, q)
        -
        \LambdaOp_\theta(\tmu_m, q)
    \big\|
    \;\le\;
    L_{\text{net}} \,
    \Wass{1}(\muN,\tmu_m),
\end{equation}
where
\begin{equation}
    L_{\text{net}}
    \;\coloneqq\;
    \prod_{l=1}^L L_{\text{layer}}^{(l)}.
\end{equation}
\end{theorem}

The proof (Appendix~\ref{app:proof_bound}) simply composes the Lipschitz bound in Lemma~\ref{lemma:lipschitz} with the Lipschitz continuity of MLP and normalization layers.  
The constant $L_{\text{net}}$ can in principle grow with depth, but techniques such as spectral normalization or Lipschitz regularization  
\citep[e.g.,][]{chizat2018global,castin2024how} can be used to keep it controlled in practice.  
Crucially, Theorem~\ref{thm:bound} decouples the \emph{geometric} difficulty of compression (how small can $\Wass{1}(\muN,\tmu_m)$ be?)  
from the \emph{architectural} sensitivity of the network (how large is $L_{\text{net}}$?).

\subsection{Optimal Scaling Laws via Quantization}
\label{sec:theory_scaling}

The remaining question is: how small can $\Wass{1}(\muN,\tmu_m)$ be for a given KV budget $m$?  
This is exactly the object of classical quantization theory \citep{graf2000foundations,villani2009optimal}.

\begin{corollary}[Minimax-Optimal Scaling for KV Compression]
\label{cor:scaling}
Suppose the support of $\muN$ lies on a $d$-dimensional submanifold of $\R^D$ with mild regularity \citep{graf2000foundations}.  
Then there exists a compression operator $\Compress$ that produces a measure $\tmu_m$ with at most $m$ atoms such that
\begin{equation}
    \Wass{1}(\muN,\tmu_m)
    \;\le\;
    C \, m^{-1/d},
\end{equation}
for some constant $C$ depending on the data distribution but not on $m$.  
Moreover, for general measures supported on $d$-dimensional manifolds, this rate $m^{-1/d}$ is minimax-optimal and cannot be improved in general.  
Combined with Theorem~\ref{thm:bound}, the ICL prediction error under KV compression decays as
\begin{equation}
    \big\|
        \LambdaOp_\theta(\muN, q)
        -
        \LambdaOp_\theta(\tmu_m, q)
    \big\|
    \;=\;
    \tilde{\mathcal{O}}(m^{-1/d}),
\end{equation}
where the $\tilde{\mathcal{O}}$ hides logarithmic factors and network-dependent Lipschitz constants.
\end{corollary}

Corollary~\ref{cor:scaling} formalizes the \emph{intrinsic-dimension bottleneck} for KV compression:  
contexts whose token features lie on low-dimensional manifolds (e.g., structured scientific or medical data) are highly compressible,  
whereas high-entropy text with large intrinsic dimension is fundamentally harder to summarize, regardless of algorithmic ingenuity.

This theory directly motivates our Measure-Aware Compression (MAC) algorithms in Section~\ref{sec:method},  
which approximate the optimal quantizers predicted by \citet{graf2000foundations} using weighted K-Means and Sinkhorn-regularized OT \citep{cuturi2013sinkhorn},  
and empirically achieve the predicted $m^{-1/d}$ scaling on synthetic tasks while improving memory--accuracy tradeoffs on real-world benchmarks.

% =================================================================
% 5. MEASURE-AWARE COMPRESSION (MAC)
% =================================================================
\section{Measure-Aware Compression (MAC)}
\label{sec:method}

The theory in Section~\ref{sec:theory} suggests that an optimal KV compressor should
\emph{(i)} minimize the Wasserstein-1 distance $\Wass{1}(\muN,\tilde\mu_m)$ between the
full and compressed context measures, and \emph{(ii)} exploit the intrinsic low-dimensional
structure of token features to achieve the minimax $m^{-1/d}$ scaling. In practice, however,
solving the exact $W_1$-optimal quantization problem is computationally prohibitive for
large $N$.

We therefore introduce \emph{Measure-Aware Compression} (MAC), a family of algorithms that
approximate the Wasserstein minimizer under realistic compute budgets:

\begin{itemize}[leftmargin=*,noitemsep,topsep=2pt]
    \item \textbf{MAC-I} (Inference-only): a post-hoc KV compressor based on weighted
    K-Means, applied once during the prefill phase of generation.
    \item \textbf{MAC-T} (Training-time): an end-to-end learnable compressor that combines
    induced set attention \citep{lee2019set} with a Sinkhorn-regularized OT loss
    \citep{cuturi2013sinkhorn}.
\end{itemize}

Both variants treat the context as a probability measure and explicitly construct a
finite-support approximation $\tilde\mu_m$ rather than scoring and dropping individual
tokens.

\subsection{MAC-I: Inference-Time Compression}
\label{sec:method_inference}

For pre-trained models where weights are frozen, we compress the KV cache after prefill,
before decoding. Let $K \in \R^{N \times D}$ and $V \in \R^{N \times D}$ denote the keys
and values for a given layer and head (for simplicity we describe the single-head case).

\paragraph{Algorithm.}
MAC-I interprets the rows of $K$ as samples from the context measure $\muN$ and seeks a
finite-support approximation with $m$ atoms. Since K-Means minimizes a discrete
$\Wass{2}^2$-type quantization error, which upper-bounds $\Wass{1}^2$ on bounded domains
\citep{villani2009optimal,graf2000foundations}, it provides a natural, scalable proxy for
Wasserstein minimization.

\begin{algorithm}[h]
   \caption{MAC-I: Inference-Time KV Compression}
   \label{alg:mac_i}
\begin{algorithmic}[1]
   \STATE {\bf Input:} Keys $K \in \R^{N \times D}$, values $V \in \R^{N \times D}$,
   budget $m$.
   \STATE {\bf (Optional) Dimensionality reduction:} Project $K$ to
   $K' \in \R^{N \times d_{\text{eff}}}$ via PCA or a learned linear map, with
   $d_{\text{eff}} \ll D$.
   \STATE {\bf K-Means clustering:} Run $m$-means on $K'$ with K-Means++ seeding
   \citep{arthur2007k} for $t$ iterations to obtain centroids
   $\{\tilde{k}_j\}_{j=1}^m$ and assignments $c(i) \in \{1,\dots,m\}$.
   \STATE {\bf Value aggregation:} For each cluster $C_j = \{i : c(i) = j\}$, define
   \begin{align*}
       \tilde{v}_j &\leftarrow \frac{1}{|C_j|} \sum_{i \in C_j} V_i, \\
       w_j &\leftarrow \frac{|C_j|}{N}.
   \end{align*}
   \STATE {\bf Output:} Compressed KV cache
   $\{(\tilde{k}_j, \tilde{v}_j, w_j)\}_{j=1}^m$, representing the sparse measure
   $\tilde\mu_m = \sum_{j=1}^m w_j \delta_{(\tilde{k}_j,\tilde{v}_j)}$.
\end{algorithmic}
\end{algorithm}

At decoding time, attention is computed against the compressed keys and values,
respecting the cluster weights $w_j$ (i.e., multiplying attention logits by $w_j$ or
equivalently repeating each centroid according to its mass). For multi-head attention,
we either (a) run MAC-I independently per head (more accurate, higher cost) or
(b) cluster in a shared low-dimensional space of concatenated keys to amortize cost
across heads; we compare these options in Appendix~\ref{app:exp_details}.

\paragraph{Complexity.}
For a single layer and head, MAC-I incurs $O(t N m d_{\text{eff}})$ time and $O(mD)$
memory during prefill. Since prefill is amortized over many decoding steps in long
generation, this overhead is competitive with token-pruning methods that recompute
scores at each step \citep{zhang2024h2o,liu2024scissorhands,li2024snapkv} and with
retrieval-based schemes that repeatedly invoke a vector index
\citep{tang2024quest,gao2024retrieval}. Once compressed, decoding operates on a KV
cache of size $m$ with the same $O(m)$ per-step complexity as other KV compression
methods.

\subsection{MAC-T: Training-Time Compression}
\label{sec:method_training}

When fine-tuning or domain-specific training is allowed (e.g., in medical or scientific
settings), we can learn a compressive representation jointly with the task. MAC-T adds a
small \emph{measure-aware bottleneck} in front of standard attention.

\paragraph{Induced Set Attention as Learnable Quantizer.}
We adopt induced set attention blocks from the Set Transformer
\citep{lee2019set}: a set of $m$ learnable inducing points
$I \in \R^{m \times D}$ attends to the full context $Z_N$ and produces a compressed set
$\{\tilde{z}_j\}_{j=1}^m$. This defines a parametric map
$\muN \mapsto \tilde\mu_m^\theta = \frac{1}{m}\sum_{j=1}^m \delta_{\tilde{z}_j}$ that
plays the role of $\Compress(\cdot)$ in our theory. We then feed $\tilde\mu_m^\theta$ to
downstream attention layers instead of $\muN$.

\paragraph{Sinkhorn-regularized OT Coverage.}
Without additional structure, inducing points may collapse onto a few modes and ignore
rare but important regions of the measure. To encourage coverage of the support of
$\muN$, we add an entropic OT regularizer between the empirical context measure and the
induced compressed measure:
\begin{equation}
    \mathcal{L}_{\text{OT}}
    \;=\;
    \text{Sinkhorn}_\varepsilon(\muN, \tilde\mu_m^\theta),
\end{equation}
where $\text{Sinkhorn}_\varepsilon$ is the Sinkhorn approximation to $\Wass{1}$ with
entropy parameter $\varepsilon$ computed via iterative matrix scaling
\citep{cuturi2013sinkhorn}. The total training objective becomes
\begin{equation}
    \mathcal{L}_{\text{total}}
    \;=\;
    \mathcal{L}_{\text{task}}
    \;+\;
    \lambda \, \mathcal{L}_{\text{OT}},
\end{equation}
with $\lambda$ controlling the strength of the geometric regularization.

This design closely matches the structure predicted by our theory: the induced points
play the role of quantization centers, and the Sinkhorn loss encourages the learned
compressor to minimize a differentiable approximation of $\Wass{1}(\muN,\tilde\mu_m)$.
In contrast to retrieval-based methods with non-differentiable selection
\citep{tang2024quest,gao2024retrieval}, MAC-T is fully end-to-end trainable.

\subsection{Practical Design Choices and Comparison}
\label{sec:method_practical}

We summarize the computational profile of MAC against representative baselines below,
where $N$ is context length, $m$ is KV budget, and $D$ is hidden dimension:

\begin{table}[h]
\caption{Asymptotic complexity per layer during prefill and decoding.}
\label{tab:complexity}
\begin{center}
\begin{small}
\begin{sc}
\begin{tabular}{lcc}
\toprule
Method & Prefill & Decode (per step) \\
\midrule
Full Attention & $O(N^2 D)$ & $O(ND)$ \\
Token Pruning (H2O) & $O(N^2 D)$ & $O(mD)$ \\
Retrieval (Quest) & $O(ND + \log N)$ & $O(mD + \log N)$ \\
Quantized KV (KIVI) & $O(N^2 D)$ & $O(ND)$ \\
\textbf{MAC-I (Ours)} & $O(t N m D)$ & $O(mD)$ \\
\textbf{MAC-T (Ours)} & $O(N^2 D)$ & $O(mD)$ \\
\bottomrule
\end{tabular}
\end{sc}
\end{small}
\end{center}
\vspace{-8pt}
\end{table}

In MAC-I, the only additional cost over standard prefill is the $O(t N m D)$ clustering
step, which is typically dominated by the $O(N^2)$ attention cost when $m \ll N$ and
$t$ is small (we use $t \le 10$ in our experiments). In MAC-T, we pay the same asymptotic
cost as standard attention during training but replace the full context by $m$ inducing
tokens at inference time.

Conceptually, MAC differs from existing KV compression strategies in how it uses the
geometry of the context: rather than scoring and dropping tokens, it \emph{constructs}
a small set of representative atoms that approximate the full context measure in
Wasserstein geometry. Section~\ref{sec:exp} shows that this measure-aware design is
sufficient to recover the predicted $m^{-1/d}$ scaling on synthetic tasks and to
outperform strong baselines such as H2O and Quest on the LUNA16 medical benchmark.

% =================================================================
% 6. EXPERIMENTS
% =================================================================
\section{Experiments}
\label{sec:exp}

We empirically evaluate MAC on (i) a controlled synthetic regression task, where we can
directly test the $m^{-1/d}$ scaling predicted by Corollary~\ref{cor:scaling}, and
(ii) a real-world long-context medical benchmark (LUNA16) that reflects the intended
scientific use cases of our framework \citep{moor2023foundation,setio2017validation}.
Unless otherwise noted, we apply MAC-I at inference time; additional details, including
MAC-T results, are deferred to Appendix~\ref{app:exp_details}.

\subsection{Experimental Setup}
\label{sec:exp_setup}

\paragraph{Models.}
All experiments use standard Transformer encoders with multi-head attention
\citep{vaswani2017attention}. For each setting, we first train a baseline model on full
contexts (no KV compression), then freeze the weights and apply different KV compression
schemes only at inference time. This isolates the effect of compression from model
capacity. We report mean performance over multiple random seeds.

\paragraph{Baselines.}
We compare MAC against:
\begin{itemize}[leftmargin=*,noitemsep,topsep=1pt]
    \item \textbf{Full KV}: no compression (upper bound).
    \item \textbf{Random}: keep $m$ tokens chosen uniformly at random.
    \item \textbf{H2O} \citep{zhang2024h2o}: heavy-hitter token pruning based on accumulated
    attention scores.
    \item \textbf{Quest} \citep{tang2024quest}: query-aware page-level sparsification.
\end{itemize}
For fairness, all compressed methods use the same KV budget $m$ at test time.

\subsection{Synthetic Linear Regression: Testing the Scaling Law}
\label{sec:exp_synth}

\paragraph{Task.}
We follow the in-context linear regression setup used in prior ICL work
\citep[e.g.][]{garg2022what,akyurek2022learning,von2023transformers}. Inputs
$x_i \in \R^{d}$ are drawn from a mixture of Gaussians, outputs are
$y_i = w^\top x_i + \epsilon$, and the model must predict $y$ for new queries given a
context of $N$ labeled examples. Here $d$ is controlled, allowing us to probe the
intrinsic-dimension dependence in Corollary~\ref{cor:scaling}.

\paragraph{Protocol.}
We train a 4-layer Transformer encoder (no positional encodings) on full contexts of
length $N=1024$. At test time, we compress the KV cache to size $m \in \{8, 16, 32, 64\}$
using MAC-I, Random, and H2O, and measure the mean-squared error (MSE) on held-out
queries. MAC-I uses a low-dimensional projection of keys ($d_{\text{eff}} \ll D$) before
K-Means, as described in Section~\ref{sec:method_inference}.

\paragraph{Results.}
Figure~\ref{fig:synth_scaling} shows MSE versus $m$ on a log--log scale. MAC-I closely
follows a straight line with slope consistent with the theoretical
$\mathcal{O}(m^{-1/d})$ prediction, while Random pruning decays much more slowly and H2O
plateaus early. Intuitively, H2O over-focuses on tokens that are salient for the current
query but fails to preserve a good geometric summary of the underlying data distribution,
whereas MAC explicitly approximates the context measure.

\begin{figure}[t]
\vskip 0.15in
\begin{center}
\centerline{\includegraphics[width=0.95\columnwidth]{figs/example-image-a}}
\caption{\textbf{Synthetic linear regression.}
Test MSE versus KV budget $m$ (log--log scale). MAC-I tracks the predicted
$m^{-1/d}$ slope, while score-based pruning (H2O) saturates for moderate $m$.}
\label{fig:synth_scaling}
\end{center}
\vskip -0.2in
\end{figure}

\subsection{LUNA16 Pulmonary Nodule Detection}
\label{sec:exp_luna}

\paragraph{Task and Data.}
LUNA16 is a public benchmark for pulmonary nodule detection from chest CT scans
\citep{setio2017validation}. Each scan consists of $200$--$500$ slices; following
\citet{moor2023foundation}, we view a study as a long sequence of 2D slices and ask the
model to predict, at study level, whether it contains nodules. This is a prototypical
long-context medical task where only a few slices (nodules) are clinically critical, and
the rest provide background.

\paragraph{Model and Features.}
We extract per-slice visual features using a ResNet-50 backbone pretrained on ImageNet,
projected to a $D$-dimensional embedding. A 4-layer Transformer encoder with multi-head
attention processes the sequence, and a pooled representation feeds a binary classifier.
The model is trained on full sequences; at test time we apply MAC-I, H2O, Quest, or
windowed attention under the same KV budget $m$.

\paragraph{Results.}
Table~\ref{tab:luna} reports accuracy, peak VRAM usage, and relative decoding speed for
the LUNA16 test set with a fixed budget of $m=64$ cached tokens (roughly $10$–$20\%$ of
the original context length).

\begin{table}[h]
\caption{LUNA16 nodule detection with KV budget $m=64$.}
\label{tab:luna}
\begin{center}
\begin{small}
\begin{sc}
\begin{tabular}{lccc}
\toprule
Method & Accuracy & VRAM & Speedup \\
\midrule
Full KV             & 94.5\% & 24GB & 1.0$\times$ \\
Local window        & 78.2\% & \textbf{2GB} & 4.5$\times$ \\
H2O \citep{zhang2024h2o}   & 85.1\% & 4GB  & 4.2$\times$ \\
Quest \citep{tang2024quest} & 88.4\% & 3GB  & 4.3$\times$ \\
\textbf{MAC-I (ours)} & \textbf{92.3\%} & 4GB  & 4.1$\times$ \\
\bottomrule
\end{tabular}
\end{sc}
\end{small}
\end{center}
\vspace{-8pt}
\end{table}

MAC-I achieves substantially better accuracy than H2O and Quest at similar memory and
speed, closing most of the gap to the Full KV upper bound while reducing VRAM by
$\sim 6\times$. Qualitative inspection of the learned clusters shows that MAC tends to
assign dedicated centroids to rare nodule-like slices in feature space, whereas H2O
sometimes evicts them in favor of more frequently attended but less informative slices.

\subsection{Summary}
\label{sec:exp_summary}

Across synthetic and real-world settings, the empirical behavior of MAC closely matches
our measure-theoretic analysis: (i) on controlled regression tasks, MAC follows the
$m^{-1/d}$ error decay predicted by quantization theory; (ii) on LUNA16, MAC delivers
state-of-the-art memory--accuracy tradeoffs compared to strong KV compression baselines,
supporting the view that KV compression is best understood as constructing an optimal
geometric summary of the context distribution rather than dropping low-score tokens.

% =================================================================
% 7. LIMITATIONS AND DISCUSSION
% =================================================================
\section{Limitations and Discussion}
\label{sec:discussion}

Our results support the view that Transformers behave as \emph{measure learners} and that
KV compression is fundamentally a quantization problem in Wasserstein geometry. At the
same time, the framework has several limitations that point to concrete directions for
future work.

\paragraph{Computational Cost of Measure-Aware Compression.}
MAC-I trades off better geometric fidelity for additional prefill cost: K-Means-based
clustering scales as $O(t N m D)$, which is cheaper than full $O(N^2 D)$ attention but
more expensive than simple top-$k$ token scoring \citep{zhang2024h2o,liu2024scissorhands}
or scalar quantization \citep{liu2024kivi}. In the settings we study, this overhead is
amortized over long decoding horizons, but for very short sequences or extremely tight
latency constraints, lighter-weight approximations (e.g., hierarchical clustering or
coresets \citep{ge2024model}, or retrieval-style vector indices \citep{gao2024retrieval})
may be preferable. A promising direction is to co-design MAC with the underlying
inference stack so that geometric compression reuses infrastructure already deployed for
retrieval \citep{tang2024quest,bertsch2024unlimiformer}.

\paragraph{From Static Measures to Causal, Infinite Context.}
Our theory treats the context as a static measure and is explicitly permutation
invariant. In contrast, autoregressive LMs and streaming models operate with a
time-evolving context $\mu_t$ and causal masking. Extending our Wasserstein stability
results to this setting likely requires “time-lifting’’ or path-space constructions
\citep{furuya2025transformers,castin2024how}, and a careful treatment of recurrent or
infinite-context architectures such as Infini-attention
\citep{munkhdalai2024leave,bertsch2024unlimiformer} and state-space models
\citep{gu2024mamba}. Developing measure-theoretic compression schemes that remain
consistent along the entire generation trajectory is an important open problem.

\paragraph{Tightness of Lipschitz Bounds.}
The KV compression error bound in Theorem~\ref{thm:bound} is worst-case and depends on a
global network Lipschitz constant $L_{\text{net}}$, which can be loose for deep
Transformers. In practice, normalization, residual connections, and attention structure
often induce much smaller effective sensitivities than the naive product of per-layer
constants would suggest \citep{castin2024how,sander2022sinkformers,chizat2018global}.
Characterizing the \emph{effective} Lipschitz behavior of trained models and tying it to
their KV compressibility is an interesting direction, and could motivate explicit
regularization schemes that improve both robustness and compressibility.

\paragraph{Scope of Empirical Validation.}
Our experiments focus on medium-scale Transformers and two representative settings:
synthetic linear regression and long-context medical imaging
\citep{moor2023foundation,setio2017validation}. While this is sufficient to observe the
predicted $m^{-1/d}$ scaling and to outperform strong KV compression baselines, it does
not yet demonstrate benefits at the scale of frontier language models
\citep{touvron2023llama,team2024gemini,ding2024longrope}. Scaling MAC to multi-billion
parameter models and diverse domains (e.g., code, multi-modal scientific data) is
primarily an engineering challenge but may also reveal new phenomena not captured by our
current theory.

\paragraph{Interaction with Other Efficiency Techniques.}
We study MAC on top of standard softmax attention \citep{vaswani2017attention}. In
practice, KV compression is often combined with linear/streaming attention
\citep{katharopoulos2020transformers,xiao2024efficient}, quantization
\citep{liu2024kivi}, and retrieval \citep{tang2024quest,gao2024retrieval}. Understanding
how measure-aware compression composes with these techniques—e.g., whether MAC can serve
as a geometric backbone into which retrieval or quantization plug as special cases—is an
exciting direction for building unified long-context systems.

% =================================================================
% 8. CONCLUSION
% =================================================================
\section{Conclusion}
\label{sec:conclusion}

We proposed a measure-theoretic view of In-Context Learning in Transformers, modeling the
KV cache as a finite-support approximation of a context measure. Under mild regularity
assumptions, we showed that kernel attention is Lipschitz with respect to the
Wasserstein-1 distance, and combined this with classical quantization theory to derive a
minimax-optimal $m^{-1/d}$ scaling law for KV compression error. Guided by this
geometry, we introduced Measure-Aware Compression (MAC), which approximates
Wasserstein-optimal quantization via K-Means at inference time and Sinkhorn-regularized
induced attention during training. MAC empirically follows the predicted scaling on
synthetic regression and achieves favorable memory–accuracy tradeoffs on a real-world
medical benchmark compared to state-of-the-art KV compression methods.

More broadly, our results support the thesis that Transformers are not only sequence
models but \emph{measure learners}, and that KV compression is best understood as
constructing an optimal geometric summary of the context distribution rather than
dropping low-score tokens. We hope this perspective will help bridge theory and systems
for long-context models, and inspire new compression and memory architectures grounded
in optimal transport and the geometry of probability measures.

% =================================================================
% 8. BROADER IMPACT
% =================================================================
\section*{Broader Impact}
\label{sec:broader_impact}

Our work is methodological: we study Transformers as measure learners and design
geometric KV compression schemes. A positive consequence is that long-context
models may become more memory- and energy-efficient, making it easier for
resource-limited settings (e.g., hospitals or smaller labs) to use foundation
models on long medical or scientific sequences
\citep{moor2023foundation,setio2017validation,team2024gemini,ding2024longrope,gu2024mamba}.

At the same time, cheaper and more scalable long-context inference also lowers
the barrier for misuse, for example in privacy-invasive logging or more
targeted generation of harmful content. Our methods do not address dataset
biases, fairness, or privacy by themselves; they should be combined with
appropriate data governance and safety measures
\citep{vaswani2017attention,touvron2023llama}. Overall, we view our
contribution as clarifying the geometry of in-context learning; its societal
impact will depend on how it is instantiated and regulated in downstream
applications.

% =================================================================
% BIBLIOGRAPHY
% =================================================================
\bibliography{refs}
\bibliographystyle{icml2026}

% =================================================================
% APPENDIX
% =================================================================
\newpage
\appendix
\onecolumn

\section{Detailed Proofs of Theoretical Results}
\label{app:proofs}

In this appendix we provide proofs for the main theoretical results in
Section~\ref{sec:theory}. Throughout, $\OmegaSpace \subset \mathbb{R}^d$ is a compact
metric space, $\Pspace$ is the set of Borel probability measures on $\OmegaSpace$, and
$\Wass{1}$ denotes the 1-Wasserstein distance on $\Pspace$
\citep{villani2009optimal}.

\subsection{Notation and Preliminaries}

For a bounded measurable function $f:\OmegaSpace \to \mathbb{R}^k$ and
$\mu \in \Pspace$, we write
\[
    \langle f,\mu\rangle \;\coloneqq\; \int_{\OmegaSpace} f(z)\, d\mu(z).
\]
The Kantorovich--Rubinstein dual formulation of $\Wass{1}$ states that for all
$\mu,\nu \in \Pspace$,
\begin{equation}
    \Wass{1}(\mu,\nu)
    =
    \sup_{\|f\|_{\mathrm{Lip}} \le 1}
    \left|
        \langle f,\mu\rangle
        -
        \langle f,\nu\rangle
    \right|,
\end{equation}
where the supremum runs over all 1-Lipschitz real-valued functions on $\OmegaSpace$
\citep{villani2009optimal}.

Let $\mathcal{H}_K$ be the reproducing kernel Hilbert space (RKHS) associated with a
positive definite kernel $\kappa:\OmegaSpace \times \OmegaSpace \to \mathbb{R}$. The
kernel mean embedding of $\mu\in\Pspace$ is
\[
    \phi(\mu) \;\coloneqq\; \int_{\OmegaSpace} \kappa(\cdot,z)\, d\mu(z) \in \mathcal{H}_K.
\]
A kernel is \emph{characteristic} if the map $\mu\mapsto \phi(\mu)$ is injective; see
\citet{muandet2017kernel,sriperumbudur2010hilbert}.

\subsection{Proof of Theorem~\ref{thm:universality} (Universality on Measures)}
\label{app:proof_universality}

\paragraph{Restatement.}
Under Assumption~\ref{ass:regularity}, for any continuous functional
$F \in C(\Pspace \times \OmegaSpace)$ and any $\varepsilon>0$, there exists a Transformer
$\Lambda_\theta$ with kernel-attention heads and MLP layers such that
\[
    \sup_{\mu\in\Pspace,q\in\OmegaSpace}
    \big\|F(\mu,q)-\Lambda_\theta(\mu,q)\big\| < \varepsilon.
\]

\paragraph{Proof sketch.}
We follow the strategy of \citet{furuya2025transformers}, adapting it to our
kernel-attention parameterization.

\textbf{1. Compactness.}
By Assumption~\ref{ass:regularity}, $\OmegaSpace$ is compact. By Prokhorov's theorem,
$\Pspace$ is compact under the topology of weak convergence, which is metrized by
$\Wass{1}$ on bounded domains \citep{villani2009optimal}. Therefore
$\mathcal{X} \coloneqq \Pspace \times \OmegaSpace$ is compact Hausdorff.

\textbf{2. Function algebra generated by attention heads.}
Consider the class $\mathcal{A}$ of functions from $\mathcal{X}$ to $\mathbb{R}^k$
representable by a finite composition of:
(i) kernel-attention maps of the form
\[
    (\mu,q) \mapsto 
    \Attn(q;\mu) 
    =
    \frac{\int \kappa(q,z)\,\varphi(z)\, d\mu(z)}
         {\int \kappa(q,z)\, d\mu(z)}
\]
with bounded Lipschitz $\varphi$; and
(ii) finite-width MLPs with continuous activations applied to finite collections of such
features and $q$. Because MLPs are universal on compact subsets of Euclidean space,
$\mathcal{A}$ is closed under finite linear combinations and pointwise products, and
contains constant functions. Hence $\mathcal{A}$ is a subalgebra of $C(\mathcal{X})$.

\textbf{3. Separation of points.}
Let $(\mu,q)\neq(\nu,r)$ in $\mathcal{X}$. We must show there exists $f\in\mathcal{A}$
such that $f(\mu,q)\neq f(\nu,r)$.

\emph{Case 1: $q\neq r$.}  
Use an MLP acting only on the query: the map
$q \mapsto \mathrm{MLP}(q)$ is universal on compact $\OmegaSpace$, so there exists an MLP
that separates $q$ and $r$.

\emph{Case 2: $q=r$ and $\mu\neq\nu$.}  
Consider the kernel mean embeddings
\[
    \phi(\mu) = \int \kappa(\cdot,z)\,d\mu(z),
    \qquad
    \phi(\nu) = \int \kappa(\cdot,z)\,d\nu(z).
\]
Since $\kappa$ is characteristic, $\phi(\mu)\neq\phi(\nu)$
\citep{muandet2017kernel,sriperumbudur2010hilbert}. Therefore there exists some
evaluation point $q^\star\in\OmegaSpace$ such that
\[
    \int \kappa(q^\star,z)\, d\mu(z)
    \;\neq\;
    \int \kappa(q^\star,z)\, d\nu(z).
\]
This evaluation is precisely the (unnormalized) numerator of an attention head with
query $q^\star$ and value map $\varphi\equiv 1$. Thus an attention head yields a scalar
feature that separates $\mu$ and $\nu$; combining it with a simple linear readout gives a
function in $\mathcal{A}$ that separates $(\mu,q^\star)$ and $(\nu,q^\star)$.

Therefore $\mathcal{A}$ separates points in $\mathcal{X}$.

\textbf{4. Stone--Weierstrass.}
By the Stone--Weierstrass theorem for compact Hausdorff spaces, any subalgebra
$\mathcal{A}$ of $C(\mathcal{X})$ that contains constants and separates points is dense
in $C(\mathcal{X})$ under the uniform norm. Hence for any $F\in C(\mathcal{X})$ and
$\varepsilon>0$, there exists $f\in\mathcal{A}$ with
$\sup_{(\mu,q)\in\mathcal{X}}\|F(\mu,q)-f(\mu,q)\| < \varepsilon$.

Any such $f$ can be realized by a Transformer with finitely many kernel-attention heads
and MLP layers, since $\mathcal{A}$ is defined exactly as the set of functions realized
by this architecture. This concludes the proof. \hfill$\square$

\subsection{Proof of Lemma~\ref{lemma:lipschitz} (Lipschitz Continuity)}
\label{app:proof_lipschitz}

\paragraph{Restatement.}
Under Assumption~\ref{ass:regularity}, for any fixed query $q\in\OmegaSpace$, the map
$\mu\mapsto \Attn(q;\mu)$ is Lipschitz with respect to $\Wass{1}$:
\[
    \big\|\Attn(q;\mu)-\Attn(q;\nu)\big\|
    \le
    L_{\Attn}\,\Wass{1}(\mu,\nu),
    \quad \forall \mu,\nu\in\Pspace.
\]

\paragraph{Proof.}
Fix $q$ and define the scalar- or vector-valued functions
\[
    f_N(z) \coloneqq \kappa(q,z)\,\varphi(z),
    \qquad
    f_D(z) \coloneqq \kappa(q,z),
\]
so that
\[
    N(\mu) \coloneqq \langle f_N,\mu\rangle,
    \qquad
    D(\mu) \coloneqq \langle f_D,\mu\rangle,
    \qquad
    \Attn(q;\mu) = \frac{N(\mu)}{D(\mu)}.
\]

By Assumption~\ref{ass:regularity}, $\kappa$ and $\varphi$ are bounded and Lipschitz on
compact $\OmegaSpace$, hence $f_N$ and $f_D$ are bounded and Lipschitz. Let
$\|f_N\|_\infty \le B_N$, $\|f_D\|_\infty \le B_D$, and
$\mathrm{Lip}(f_N) \le L_N$, $\mathrm{Lip}(f_D) \le L_D$.

Using the dual formulation of $\Wass{1}$,
\[
    \|N(\mu)-N(\nu)\|
    =
    \left\|\int f_N\, d(\mu-\nu)\right\|
    \le
    L_N \, \Wass{1}(\mu,\nu),
\]
and similarly
\[
    |D(\mu)-D(\nu)|
    \le
    L_D \, \Wass{1}(\mu,\nu).
\]
Moreover, $|D(\mu)|\le B_D$ for all $\mu$ and, by Assumption~\ref{ass:regularity}, there
exists $\delta>0$ such that $D(\mu)\ge\delta$ for all relevant $\mu$.

We now bound the difference of the quotients:
\begin{align*}
    \Attn(q;\mu) - \Attn(q;\nu)
    &= \frac{N(\mu)}{D(\mu)} - \frac{N(\nu)}{D(\nu)} \\
    &= \frac{N(\mu)D(\nu) - N(\nu)D(\mu)}{D(\mu)D(\nu)} \\
    &= \frac{N(\mu)\big(D(\nu)-D(\mu)\big)
        + D(\mu)\big(N(\mu)-N(\nu)\big)}
       {D(\mu)D(\nu)} .
\end{align*}
Taking norms and using $|D(\mu)|,|D(\nu)|\ge\delta$,
\begin{align*}
    \big\|\Attn(q;\mu)-\Attn(q;\nu)\big\|
    &\le
    \frac{\|N(\mu)\|\,|D(\nu)-D(\mu)|
          + |D(\mu)|\,\|N(\mu)-N(\nu)\|}
         {|D(\mu)D(\nu)|} \\
    &\le
    \frac{B_N L_D + B_D L_N}{\delta^2} \, \Wass{1}(\mu,\nu).
\end{align*}
Thus the map is Lipschitz with constant
$L_{\Attn} \coloneqq (B_N L_D + B_D L_N)/\delta^2$. \hfill$\square$

\subsection{Proof of Theorem~\ref{thm:bound} (KV Compression Bound)}
\label{app:proof_bound}

\paragraph{Restatement.}
Let $\Lambda_\theta = \Lambda^{(L)}\circ\cdots\circ\Lambda^{(1)}$ be an $L$-layer
Transformer, where each layer $\Lambda^{(l)}$ is $L_{\text{layer}}^{(l)}$-Lipschitz in
its measure argument. For any empirical context measure $\muN$ and compressed measure
$\tilde\mu_m$,
\[
    \big\|\Lambda_\theta(\muN,q) - \Lambda_\theta(\tilde\mu_m,q)\big\|
    \le
    L_{\text{net}}\,\Wass{1}(\muN,\tilde\mu_m),
\]
with $L_{\text{net}} = \prod_{l=1}^L L_{\text{layer}}^{(l)}$.

\paragraph{Proof.}
Define $h^{(0)}(\mu,q) = (\mu,q)$ and for $l\ge 1$,
\[
    h^{(l)}(\mu,q)
    \;=\;
    \Lambda^{(l)}\big(h^{(l-1)}(\mu,q)\big).
\]
Then $\Lambda_\theta(\mu,q) = h^{(L)}(\mu,q)$.

Let $\mu=\muN$ and $\nu=\tilde\mu_m$. By Lipschitzness of each layer,
\[
    \big\|h^{(1)}(\mu,q)-h^{(1)}(\nu,q)\big\|
    \le
    L_{\text{layer}}^{(1)} \,\Wass{1}(\mu,\nu).
\]
Applying the same argument recursively,
\begin{align*}
    \big\|h^{(2)}(\mu,q)-h^{(2)}(\nu,q)\big\|
    &\le
    L_{\text{layer}}^{(2)} \,
    \big\|h^{(1)}(\mu,q)-h^{(1)}(\nu,q)\big\| \\
    &\le
    L_{\text{layer}}^{(2)} L_{\text{layer}}^{(1)} \,
    \Wass{1}(\mu,\nu),
\end{align*}
and, by induction,
\[
    \big\|h^{(L)}(\mu,q)-h^{(L)}(\nu,q)\big\|
    \le
    \Big(\prod_{l=1}^L L_{\text{layer}}^{(l)}\Big)\,
    \Wass{1}(\mu,\nu).
\]
Setting $\mu=\muN$ and $\nu=\tilde\mu_m$ yields the stated bound. \hfill$\square$

\subsection{Proof of Corollary~\ref{cor:scaling} (Scaling Law)}
\label{app:scaling}

\paragraph{Restatement.}
If the support of $\muN$ lies on a $d$-dimensional submanifold of $\mathbb{R}^D$ with
mild regularity, there exists a compression operator producing $\tilde\mu_m$ with at
most $m$ atoms such that
\[
    \Wass{1}(\muN,\tilde\mu_m)
    \le
    C \, m^{-1/d},
\]
and this rate is minimax-optimal in general. Combining with
Theorem~\ref{thm:bound} yields an ICL prediction error of order
$\tilde{\mathcal{O}}(m^{-1/d})$.

\paragraph{Proof sketch.}
The claim is a direct application of classical quantization theory for probability
measures \citep{graf2000foundations,villani2009optimal}. For completeness, we briefly
outline the relevant result.

Let $\mathcal{Q}_m$ be the set of all probability measures with support size at most $m$.
Define the optimal $L^1$ quantization error of $\mu$ at level $m$ as
\[
    e_1(\mu,m)
    \;\coloneqq\;
    \inf_{\nu\in\mathcal{Q}_m} \Wass{1}(\mu,\nu).
\]
Under mild assumptions on $\mu$—in particular, that it admits a density with respect to
the $d$-dimensional Hausdorff measure on a compact $d$-dimensional manifold—there exist
constants $0<c\le C<\infty$ such that
\[
    c\,m^{-1/d}
    \le
    e_1(\mu,m)
    \le
    C\,m^{-1/d}
    \quad\text{for all sufficiently large }m.
\]
The upper bound is achieved by appropriately chosen $m$-point quantizers; the lower bound
shows that no other scheme can improve the rate uniformly over a rich class of measures
\citep{graf2000foundations}.

Applying this to $\mu=\muN$ yields the claimed rate
$\Wass{1}(\muN,\tilde\mu_m) \le C m^{-1/d}$ for some (possibly data-dependent) $C$.
Plugging into Theorem~\ref{thm:bound} gives
\[
    \big\|
        \Lambda_\theta(\muN,q)
        -
        \Lambda_\theta(\tilde\mu_m,q)
    \big\|
    \le
    L_{\text{net}}\,C\,m^{-1/d},
\]
which we write as $\tilde{\mathcal{O}}(m^{-1/d})$ to emphasize that $L_{\text{net}}$ and
logarithmic factors are absorbed into the constant. \hfill$\square$

% =================================================================
% B. EXPERIMENTAL DETAILS AND ADDITIONAL RESULTS
% =================================================================
\section{Experimental Details and Additional Results}
\label{app:exp_details}

We provide additional details on the experimental setups in Section~\ref{sec:exp}, as
well as ablations and implementation notes.

\subsection{Synthetic Linear Regression}
\label{app:exp_synth_details}

\paragraph{Data generation.}
We consider input dimension $d\in\{10,20\}$ to probe the dependence on intrinsic
dimension. For each task instance, we sample a weight vector $w\sim\mathcal{U}(-1,1)^d$
and generate a context of $N=1024$ pairs $(x_i,y_i)$:
\begin{itemize}[leftmargin=*,noitemsep,topsep=1pt]
    \item $x_i$ is drawn from a mixture of three Gaussians on $\mathbb{R}^d$, with
    randomly sampled means and shared covariance $0.1^2 I$. This yields a non-uniform,
    low-dimensional data manifold.
    \item $y_i = w^\top x_i + \epsilon_i$ with
    $\epsilon_i \sim \mathcal{N}(0,\sigma^2)$, $\sigma=0.05$.
\end{itemize}
We generate fresh task instances at test time to avoid overfitting to a fixed $w$.

\paragraph{Model architecture.}
We use a 4-layer Transformer encoder with hidden dimension $D=128$, 4 attention heads,
and GELU activations. No positional encodings are used so that the model is encouraged
to treat the context as a set \citep{lee2019set}. We prepend a dedicated query token and
train the model to predict its target value given the rest of the context.

\paragraph{Training and evaluation.}
The model is trained with Adam for 10\,000 steps, batch size 64, learning rate $10^{-3}$,
using full KV cache. At evaluation time, we freeze weights and apply KV compression with
different budgets $m$. For MAC-I, we project keys to $d_{\text{eff}}=16$ dimensions using
a learned linear map, then run Mini-Batch K-Means with K-Means++ initialization
\citep{arthur2007k}. We average results over 5 seeds.

\subsection{LUNA16 Pulmonary Nodule Detection}
\label{app:exp_luna_details}

\paragraph{Preprocessing.}
We follow the LUNA16 protocol \citep{setio2017validation}. CT volumes are resampled to a
fixed voxel spacing, clipped to a standardized Hounsfield range, and normalized. We then
extract 2D axial slices and center-crop each slice to a fixed resolution. Each study
typically yields $N\approx 200$--$500$ slices.

\paragraph{Feature extractor and Transformer.}
We use a ResNet-50 backbone pretrained on ImageNet to extract per-slice feature vectors
from the penultimate layer ($2048$-dim). A linear projection maps these to a Transformer
hidden dimension $D=256$. We stack a 4-layer Transformer encoder with 8 heads and apply
max-pooling over the sequence followed by a 2-layer MLP for binary classification.

\paragraph{Training.}
We train on full sequences using cross-entropy loss, with AdamW optimizer, learning rate
$10^{-4}$, batch size 16, and cosine learning rate schedule. We use early stopping based
on validation accuracy.

\paragraph{MAC-I and baselines.}
At test time:
\begin{itemize}[leftmargin=*,noitemsep,topsep=1pt]
    \item \textbf{MAC-I}: we apply PCA to reduce key dimension to $d_{\text{eff}}=32$
    and run Mini-Batch K-Means with $m=64$ clusters and $t=10$ iterations. Cluster
    centroids and masses are used as compressed KV states.
    \item \textbf{H2O}: we implement heavy-hitter token selection as in
    \citet{zhang2024h2o}, accumulating attention scores over layers and pruning to the
    top-$m$ tokens.
    \item \textbf{Quest}: we divide the sequence into pages of 16 slices and apply
    query-aware page selection following \citet{tang2024quest}.
    \item \textbf{Local window}: we restrict attention to a symmetric local window of
    size $m$ around the query.
\end{itemize}
We report accuracy, peak GPU memory, and per-step decoding time averaged over the test
set. All methods use the same underlying model and differ only in KV handling.

\subsection{MAC-T Implementation Details}
\label{app:exp_mact_details}

\paragraph{Induced set attention.}
For MAC-T experiments (not shown in the main text due to space), we insert an induced
set attention block \citep{lee2019set} between the feature extractor and the first
Transformer layer. The block uses $m=64$ inducing tokens $I\in\mathbb{R}^{m\times D}$ and
produces a compressed set $\{\tilde{z}_j\}_{j=1}^m$ that replaces the original sequence.

\paragraph{Sinkhorn loss.}
We approximate the OT distance between the empirical measure over features
$\muN = \frac{1}{N}\sum_{i=1}^N\delta_{z_i}$ and the compressed measure
$\tilde\mu_m^\theta = \frac{1}{m}\sum_{j=1}^m\delta_{\tilde{z}_j}$ using the entropic
Sinkhorn divergence \citep{cuturi2013sinkhorn} with regularization
$\varepsilon=0.05$ and 10 iterations. The OT term is added to the task loss with weight
$\lambda=0.1$.

\paragraph{Observed behavior.}
MAC-T typically learns inducing tokens that cover different semantic regions of feature
space (e.g., normal lung tissue, nodule-like patterns), and achieves slightly better
accuracy than MAC-I at the same budget $m$, at the cost of additional training
overhead. The qualitative behavior is consistent with our interpretation of the inducing
tokens as learned quantization centers guided by OT geometry.

\subsection{Additional Ablations}

Due to space, we summarize key trends:

\begin{itemize}[leftmargin=*,noitemsep,topsep=1pt]
    \item \textbf{Effect of intrinsic dimension.} On synthetic regression, increasing the
    input dimension $d$ steepens the slope of the error curve in log--log space for all
    methods, but MAC-I remains closest to the theoretical $m^{-1/d}$ prediction, while
    Random and H2O deviate more strongly as $d$ grows.
    \item \textbf{Budget sensitivity.} On LUNA16, MAC-I maintains a smooth accuracy
    degradation as $m$ decreases, degrading gracefully down to $m=32$. H2O and Quest
    exhibit sharper drops once $m$ falls below $\approx 10\%$ of the original context.
    \item \textbf{Head-wise vs. shared clustering.} For moderate budgets, clustering in a
    shared key space across heads achieves nearly the same accuracy as per-head
    clustering, while being cheaper. At very small budgets, per-head clustering yields a
    modest but consistent improvement, suggesting some head-specific geometry.
\end{itemize}

These results further support the view that explicitly respecting the measure structure
of the context---instead of operating on tokens independently---is crucial for robust KV
compression in both synthetic and real-world long-context settings.


\end{document}